\addcontentsline{toc}{chapter}{Agradecimentos}

Gostaria de expressar minha sincera gratidão ao meu orientador, Dr. Julio Camargo, por acreditar neste projeto desde o início. Sua confiança em uma proposta ambiciosa, mesmo diante das limitações impostas pela minha rotina profissional, foi fundamental para que este trabalho se tornasse realidade.

Agradeço também aos integrantes do Grupo do Rio pela colaboração e disposição em ajudar ao longo desta pesquisa. Um agradecimento especial aos meus ex-orientadores, Dr. Marcelo Assafin e Dr. Sérgio Santos-Filho, cuja orientação durante a Iniciação Científica construiu a base sólida que tornou esta etapa muito mais leve.

Registro aqui minha gratidão ao Observatório Nacional por acolher meu projeto, mesmo em condições atípicas. Ser aceito como mestrando, ainda que sem vínculo de bolsa e conciliando atividades profissionais, foi uma oportunidade de ouro.

Aos meus amigos e à minha família, obrigado por estarem presentes durante toda essa jornada. Em especial, aos meus pais, Cesar Vidal Cunha e Adriana Vieira Laidler, e ao meu irmão, Matheus Laidler Vidal Cunha, por todo o amor, incentivo e apoio nos momentos mais difíceis. Agradeço também ao apoio da minha namorada, Daniela Ferreira de Monte Lima, que acompanhou meus devaneios sobre trabalhar com astrofísica desde a adolescência e esteve ao meu lado quando terminei esta dissertação.

Aos meus tios, Cláudio Vidal e Cláudia Lassance, agradeço pelo apoio constante e por terem me proporcionado, desde cedo, o aprendizado da língua inglesa, ferramenta indispensável na minha formação. Às minhas tias Alessandra Vieira Laidler e Christiane Vieira Laidler, pelo carinho e incentivo. Aos meus primos Yan Lassance e Raphael Lassance, por todo o apoio ao longo do caminho.

Este trabalho é dedicado a todos vocês.
